\input ../preamble
\def\@cntrprefix{\the\counter{section}.}

\centerline{\setfontandscale{bf}{20pt}Schubert Calculus}
\smallskip
\centerline{\it Ari Feiglin}

\section{Grassmanians}

\bnote

    Subvarieties are assumed to be closed.

\enote

The focus of this discussion is on Schubert calculus, in particular the sort of problems it solves in
enumerative geometry.
In general the problem is to determine the number of spaces (of a given dimension) intersect some set
of given spaces.
We are dealing here with linear spaces and thus it is natural to consider a form of topology on linear
subspaces.

\bdefn

    Let $V$ be an $n$-dimensional vector space, then for $k\leq n$ we define $V$'s {\emphcolor Grassmanian}
    to be the space
    $$ G(k,V) = \set{\Lambda\leq V}[\hbox{$\Lambda$ is $k$-dimensional}] . $$

\edefn

The Grassmanian will of course be given the structure of a variety, in due time.
First of all, if $\Psi\colon V\cong V'$ is an isomorphism of vector spaces then their Grassmanians are in
bijection
$$ G(k,V) \lisoto G(k,V'),\qquad \Lambda\mapsto\Psi\Lambda , $$
and this is an isomorphism of varieties.
Writing
$$ G(k,n) = G(k,\bF^n) , $$
we see that for $n$-dimensional $V$, $G(k,V)\cong G(k,n)$.

Note further that there is an isomorphism
$$ G(k,V) \lisoto G(n-k,V^*),\qquad \Lambda\mapsto\ann\Lambda = \set{\phi\in V^*}[\phi\Lambda=0] , $$
which is again an isomorphism of varieties.
Thus we have $G(k,n)\cong G(n-k,n)$.

Now we proceed with endowing $G(k,V)$ a variety structure.
For $\Lambda\in G(k,V)$ consider the exterior power
$$ \bigwedge^k\Lambda \leq \bigwedge^kV , $$
which is $1$-dimensional.
Thus there is a natural map
$$ G(k,V) \to \bP\parens{\bigwedge^kV} \cong \bP^{\binom nk-1},\qquad \Lambda \mapsto \bigwedge^k\Lambda . $$
This map is called the {\emphcolor Pl\"ucker embedding}.
We will continue to use $N=\binom nk-1$.

\blemm

    The Pl\"ucker embedding is indeed injective, and its image in $\bP(V^k)$ is a subvariety.

\elemm

\bproof

    Choose basis vectors $v_1,\dots,v_k\in\Lambda$ so that $\eta=v_1\wedge\cdots\wedge v_k$ spans
    $\Lambda^k$.
    Considering the exterior algebra $\bigwedge V$, a vector $v\in V$ annihiltes $\eta$ (meaning
    $v\wedge\eta$) iff $v$ is in the span of $v_1,\dots,v_k$, i.e.\ $v\in\Lambda$.
    Thus we can determine $\Lambda$ from its image in $\bP(V^k)$ by
    $$ \Lambda = \set{v\in V}[v\wedge\eta=0] , $$
    so the Pl\"ucker embedding is injective.

    Extend a basis $v_1,\dots,v_k\in\Lambda$ to a basis of $V$.
    Then we can write a $\eta\in\bigwedge^kV$ as $\eta=\sum_I\eta_Iv_I$ for scalars $\eta_I$.
    Then

    Now consider the map
    $$ V \xtos{\wedge\eta} \bigwedge^{k+1}V , $$
    we claim that $\eta=v_1\wedge\cdots\wedge v_k$ for linearly independent $v_i\in V$ iff the kernel
    of this map has dimension at least $k$.
    If $\eta$ is of this form then $v_1,\dots,v_k$ are in the kernel.
    Otherwise suppose $v_1,\dots,v_k$ are linearly independent vectors in the kernel.
    We may extend these $v_i$ to a basis of $V$, so that there are scalars $\eta_I$ such that
    $$ \eta = \sum_I\eta_Iv_I,\qquad v_I = v_{i_1}\wedge\cdots\wedge v_{i_k}\hbox{ for $I=(i_1<\cdots<i_k)$} . $$
    Then we know
    $$ 0 = v_1\wedge\cdots\wedge v_k\wedge\eta = \sum_I\eta_Iv_1\wedge\cdots\wedge v_k\wedge v_I
    = \sum_{I\cap[k]=\emptyset}\eta_Iv_I , $$
    so that for any multi-index $I$ which is not $(1,\dots,k)$, we have $\eta_I=0$.
    Thus $\eta$ is a scalar multiple of $v_1\wedge\cdots\wedge v_k$, as desired.
    And thus the image of the Grassmanian is
    $$ G(k,V) = \set{\eta\in\bigwedge^kV}[\rank(V\xtos{\wedge\eta}\bigwedge^{k+1}V)\leq n-k] , $$
    which is an algebraic subset of $\bP(V^k)$ as it is the zeroset of all of the $(n-k+1)$st minors of
    the map $\bul\wedge\eta$.

    To show that $G(k,V)$ is irreducible, we note that there is a transitive action of $\gl(V)$ on
    $G(k,V)$.
    Choosing any $\Lambda\in G(k,V)$ we have that $\gl(V)\Lambda=G(k,V)$ is irreducible because
    it is the image of the irreducible $\gl(V)$ (it is an open subset of the irreducible $\bA^n$).
    \qed

\eproof

Let $e_1,\dots,e_n$ be a basis for $V$ and identify $V\cong\bF^n$ so that
$\Lambda$ is the row space of some matrix $k\times n$ matrix $A=(a_{ij})$.
Then $\bigwedge^kV$ is spanned by $e_I=e_{i_1}\wedge\cdots\wedge e_{i_k}$ for $I=(i_1<\cdots<i_k)$.
So for the basis $v_1,\dots,v_k\in\Lambda$ determined by $A$ there are scalars $p_I$ such that
$$ v_1\wedge\cdots\wedge v_k = \sum_Ip_Ie_I . $$
Then we see that $p_I$ is the determinant of the submatrix of $A$ whose columns are determined by $I$.
These are the {\emphcolor Pl\"ucker coordinates} of $A$.

\bexam

    Note that the Pl\"ucker embedding for $k=1$ gives an isomorphism $G(1,V)\cong\bP(V)$.
    So the most basic example of a Grassmanian is projective space.

\eexam

Because $G(k,V)$ is a variety it has an affine cover.
For an $(n-k)$-dimensional subspace $\Gamma\leq V$ define
$$ U_\Gamma = \set{\Lambda\in G}[\Lambda\cap\Gamma=0] $$
to be the set of all subspaces which do not intersect $\Gamma$.
This is a Zariski open subset of $G$, because taking $w_1,\dots,w_{n-k}$ a basis for $\Gamma$ and
setting $\eta=w_1\wedge\cdots\wedge w_{n-k}$, we see that
$$ U_\Gamma = \set{[\omega]\in G\subseteq\bP(V^k)}[\omega\wedge\eta\neq0] . $$
We now claim that $U_\Gamma$ is isomorphic to $\bA^{k(n-k)}$.
Let $\Omega\in U_\Gamma$ be a complementary $k$-space to $\Gamma$, so that $V=\Gamma\oplus\Omega$,
and for any $\Lambda\in U_\Gamma$ we have projections
$$ \Lambda \xtos{\pi_\Gamma} \Gamma,\qquad \Lambda \xtos{\pi_\Omega} \Omega $$
modulo $\Omega$ and $\Gamma$ respectively.
Because the kernel of $\pi_\Omega$ is $\Gamma$ we see that $\pi_\Omega$ is an isomorphism, and so
$\Lambda$ is the graph of
$$ \phi \colon \Omega \xtos{\pi_\Omega^{-1}} \Lambda \inj V \xtos{\pi_\Gamma} \Gamma . $$
Indeed its graph is
$$ \graph\phi = \set{\omega+\pi_\Gamma\pi_\Omega^{-1}\omega}[\omega\in\Omega] , $$
and any $v\in\Omega$ is in this graph via $\omega=\pi_\Omega(v)$, and they have the same dimension.

Conversely if $\phi\colon\Omega\to\Gamma$ is linear then its graph is complementary to $\Gamma$.
So we have a correspondence
$$ U_\Gamma \xtof{\Lambda\mapsto\pi_\Gamma\circ\pi_\Omega^{-1}}[\phi\mapsto\graph\phi] \hom(\Omega,\Gamma) . $$
And considering them in Pl\"ucker coordinates, we see that these are morphisms of varieties.
Because $\hom(\Omega,\Gamma)\cong\bA^{k(n-k)}$, we see that $G(k,V)$ is a variety of dimension $k(n-k)$.

\bprop

    The Grassmanian $G(k,V)$ is a projective variety of dimension $k(n-k)$.

\eprop

\bnote

    In S.\ L.\ Kleiman and Dan Laksov, they define the Grassmanian $G(k,n)$ as the collection of $k$-planes in
    $\bP^n$.
    This corresponds to our $G(k+1,n+1)$.
    We will denote this by
    $$ \bG(k,V) = \bG(k,\bP(V)) \cong G(k+1,n+1) . $$

\enote

\section{The Chow Group}

\bdefn

    Consider a complex flag $\V$ in $V$, i.e.\ a sequence of subspaces
    $$ 0 < V_1 < \cdots < V_{n-1} < V_n = V,\qquad \dim V_i = i . $$
    Furthermore suppose we have a sequence $a$
    $$ n-k \geq a_1 \geq \cdots \geq a_k \geq 0 . $$
    Define the {\emphcolor Schubert cycle} to be the closed subset
    $$ \Sigma_a(\V) = \set{\Lambda\in G(k,n)}[\dim(V_{n-k+i-a_i}\cap\Lambda)\geq i\hbox{ for all $i$}] . $$

\edefn

\bremark

    In general for a subspace $\Gamma\leq V$, the set
    $$ \set{\Lambda\in G(k,n)}[\dim(\Gamma\cap\Lambda)\geq i] $$
    is closed, as $\dim(\Gamma\cap\Lambda)=\dim\Gamma+\dim\Lambda-\dim(\Gamma+\Lambda)$ so the condition
    is equivalent to $\dim(\Gamma+\Lambda)\leq\dim\Gamma+k-i$.
    Choosing bases for $\Lambda,\Gamma$, this is iff the rank of the matrix obtained by concatenating the
    bases has rank $\leq\dim\Gamma+k-i$, which is iff all submatrices of larger size have zero determinant.

\eremark

For convenience, we write
$$ \Sigma_{a_1,\dots,a_s}(\V) = \Sigma_{a_1,\dots,a_s,0,\dots,0} $$
and
$$ \Sigma_{a^s} = \Sigma_{a,\dots,a} $$
to duplicate $a$ $s$ times.

\bexam

    Note that for any flag $\V$,
    $$ \Sigma_0(\V) = G(k,n) . $$
    That is, any $k$-dimensional subspace $\Lambda$ satisfies $\dim(V_{n-k+i}\cap\Lambda)\geq i$ for all $i$.
    Indeed,
    $$ \dim(V_{n-k+i}\cap\Lambda) = \dim V_{n-k+i} + \dim\Lambda - \dim(V_{n-k+i}+\Lambda) \geq
    n - k + i + k - n = i . $$

\eexam

\bexam

    Suppose $V_\ell$ is a subspace of dimension $\ell$, then extend it to an arbitrary complete flag $\V$
    and we see that
    $$ \Sigma_{n-k+1-\ell}(\V) = \set{\Lambda}[\Lambda\cap V_\ell\neq0] . $$
    Indeed for $i=1$ the Schubert condition gives
    $$ \dim(V_\ell\cap\Lambda) \geq 1 \iff V_\ell\cap\Lambda\neq 0 . $$
    The conditions on $i\geq1$ are tautological as shown by the previous example.

\eexam

\bexam

    Given an $\ell$-subspace $V_\ell$ then we have (for any extension $\V$)
    $$ \Sigma_{(n-\ell)^k}(\V) = \set{\Lambda}[\Lambda\subseteq V_\ell] , $$
    because the Schubert conditions gives for all $i$
    $$ \dim(\Lambda\cap V_{\ell-k+i}) \geq i , $$
    so for $i=k$ this means $\Lambda\subseteq V_\ell$, and then considering the first example the rest of
    the conditions follow.

    And similarly
    $$ \Sigma_{(n-k)^\ell}(\V) = \set{\Lambda}[V_\ell\subseteq\Lambda] $$
    because
    $$ \dim(\Lambda\cap V_i) \geq i \hbox{ for $i\leq\ell$} \iff V_\ell \leq \Lambda $$
    and the other conditions are tautological again.

\eexam

Note that if we define for sequences $a\leq b$ iff $a_i\leq b_i$ for all $i$, then
$$ \Sigma_a(\V) \subseteq \Sigma_b(\V) \iff a \leq b . $$
For a sequence $a$ define
$$ \abs a = \sum_i a_i . $$

\bthrm

    The Schubert cycle $\Sigma_a(\V)$ is an irreducible subvariety of $G(k,V)$ of codimension $\abs a$.

\ethrm

\bprop[name={intwithlin}]

    Let $e_1,\dots,e_n$ be a basis for $V$ and consider the standard flag $V_i=\gen{e_1,\dots,e_i}$.
    Let $a$ be a decreasing sequence and $b$ be the strictly increasing sequence $b_i=n-k+i-a_i$.
    Then $\Sigma_a(\V)\subseteq G(k,n)\subseteq\bP^N$ is the intersection
    $$ \Sigma_a(\V) = G(k,V) \cap \bigcap_{I\not\leq b}V(p_I) , $$
    where $p_I$ is the Pl\"ucker coordinate associated to $I=(i_1<\cdots<i_k)$.

\eprop

\bproof

    For $\Lambda\in G(k,V)$ write $\Lambda$ as the row space as the matrix $A=(a_{ij})$.
    The Pl\"ucker coordinate $p_I(\Lambda)$ is then the determinant of the submatrix of $A$ whose columns
    are given by $I$.

    Now suppose that $\Lambda\in\Sigma_a(V)$ so that $\dim(\Lambda\cap V_{b_i})\geq i$ for all $i$.
    For $I\not\leq b$ there is a $j$ such that $i_j>b_j$, and since $\dim(\Lambda\cap V_{b_j})\geq j$,
    there are $j$ linearly independent vectors in $\Lambda$ whose coordinates past $b_j$ are zero.
    If we take these vectors to be $v_1,\dots,v_j$ and extend them to a basis $v_1,\dots,v_k$ of $\Lambda$,
    then for every index in $I$ past $i_j$ (inclusive), the coordinates of $v_1,\dots,v_j$ are zero.
    Thus the submatrix of $A$ determined by $I$ has a $j\times(k-j+1)$ block of zeroes.
    That is, it is of the form

    \medskip
    $$ \left(\vcenter{\halign{\strut\hfil$#$\hfil&\@gobble{#1}\kern.2cm\vrule\kern.2cm&\hfil$#$\hfil&\rlap{\kern.5cm$#$}\cr
        \omit\hfil\rmht{\ $j-1$\ }\hfil&\omit&\omit\hfil\rmht{\ $k-j+1$\ }\hfil\cr
        *&&0&j\cr\noalign{\hrule}
        *&&*&k-j\cr
    }}\right)\kern1cm . $$
    Thus its rank is at most $(k-j)+(j-1)=k-1$ which means it is singular, so $p_I(\Lambda)=0$.

    Now suppose that $p_I(\Lambda)=0$ for all $I\not\leq b$, but $\dim(\Lambda\cap V_{b_j})<j$ for some $j$.
    So if consider the submatrix of $A$ determined by columns of index $>b_j$, it must have rank $>k-j$.
    This is because there are $<j$ rows which are also in $\Lambda\cap V_{b_j}$:

    \medskip
    $$ \left(\vcenter{\halign{\strut\hfil$#$\hfil&\@gobble{#1}\kern.2cm\vrule\kern.2cm&\hfil$#$\hfil&\rlap{\kern.5cm$#$}\cr
        \omit\hfil\rmht{\ $b_j$\ }\hfil\cr
        *&&\ 0\ &<j\cr\noalign{\hrule}
        *&&*&>k-j\cr
    }}\right)\kern1cm . $$

    Performing row operations gives us a block-diagonal matrix of full rank, and so the right column of blocks
    is linearly independent.
    So there are at least $k-j+1$ columns with indices greater than $b_j$ which are linearly independent.
    Extend these to $k$ linearly independent columns (because $A$ is of rank $k$) and let $I$ be the
    indices of these columns.
    Then the submatrix determined by $I$ is nonsingular so $p_I(\Lambda)\neq0$.
    Because at least $k-j+1$ of these indices is greater than $b_j$, and as they are ascending this means
    that $i_j>b_j$, a contradiction.
    \qed

\eproof

The reason for wanting to consider Schubert cycles follows from the above examples.
Consider the following problem.

\bprob[name={interfourq}]

    If $V_1,V_2,V_3,V_4\subseteq\bP^{2n+1}$ are four general $n$-planes, how many lines $L\subseteq\bP^{2n+1}$
    intersect all four?

\eprob

We know by the above example that $\Sigma_n(V_i)\subseteq\bG(1,2n+1)$ is the number of lines meeting $V_i$,
so we want to count the intersection
$$ Q = \bigcap\Sigma_n(V_i) . $$
Unfortunately we do not yet have the tools to do this.
%Now consider \refmath[proposition]{intwithlin}: we have $b=(n+1,2n+2)$ which means the only $I$s which are
%not bound by $b$ are of the form $I=(i_1,i_2)$ with $i_1>n+1$.
%So we choose the indices from $(n+1,2n+2]$ which gives us $\binom{n+1}2$ choices for $I$.
%Thus $\Sigma_n(V_i)=\bG(1,2n+1)\cap W_i$ where $W_i$ is an $\parens{N-\binom{n+1}2}$-plane of $\bP^N$.
%Here we have
%$$ N = \binom{2n+2}2 - 1 = 2n^2 + 3n,\qquad \binom{n+1}2 = \frac{n^2+n}2 . $$
%Now the number of 

\section{The Chow Groups}

\bdefn

    Let $X$ be any algebraic variety.
    A {\emphcolor $k$-cycle} on $X$ is a formal sum of subvarieties (closed integral subschemes) of $X$
    of dimension $k$.
    In other words the group of $k$-cycles on $X$ is the free Abelian group on $X$'s $k$-dimensional
    subvarieties, and is denoted $Z_k(X)$.
    $X$'s {\emphcolor group of cycles} is the graded Abelian group $Z(X)=\bigoplus_kZ_k(X)$.

    A cycle $\sum n_iY_i\in Z(X)$ is {\emphcolor effective} if all of the $n_i$ are nonnegative.
    A {\emphcolor divisor} is an $(n-1)$-cycle (where $n=\dim X$).

\edefn

\bdefn

    Let $K(X)$ denote the field of rational functions of $X$.


\edefn

\bye

